\subsection{Configuration Management}

The project uses GitHub for source code management and Jira for task management.

\textbf{Working Process}

\textbf{Step 1. Clone the Repository}

Each team member clones the GitHub repository to their local machine.

\begin{quote}
\textit{git clone \url{https://github.com/YaFhun06/CNPM.git}}
\end{quote}

\textbf{Step 2. Update the Jira Status}

Select the assigned task and change its status from:

\begin{center}
\textit{To Do $\rightarrow$ In Progress}
\end{center}

\textbf{Step 3. Create a Working Branch}

Team members must not work directly on the main branch.

Create a feature branch following the team's naming convention.

Branch Naming Convention: \textit{feature/sprint1-member-name}

Example:

\textit{feature/sprint1-phung}

\textit{feature/sprint2-bao}

\textbf{Step 4. Implement the Assigned Tasks}

Each member works on the assigned tasks within the project structure using Visual Studio Code.

\textbf{Step 5. Commit Changes}

Commit messages should reference the corresponding Jira Issue Key.

Commit Convention: \textit{<Jira Issue Key>: <Action> <Description>}

Where:

\begin{itemize}
  \item Jira Issue Key: CNPM-3, CNPM-15, \ldots
  \item Action: Add, Update, Fix, Refactor, Remove, Complete
  \item Description: Brief description of the changes
\end{itemize}

Example: \textit{git commit -m "CNPM-3: Add use case list"}

\textbf{Step 6. Push the Working Branch}

\textit{git push origin feature/sprint1-phung}

\textbf{Step 7. Update the Jira Status}

After the task has been completed and confirmed by the team, change its status from:

\begin{center}
\textit{In Progress $\rightarrow$ Done}
\end{center}
